% !TEX TS-program = XeLaTex
% !TEX encoding = UTF-8 Unicode

\documentclass[UTF8,AutoFakeBold,AutoFakeSlant]{beamer}
\usepackage{ctex, hyperref}
\usepackage[T1]{fontenc}

\usepackage[backend=biber,style=numeric-comp,sorting=none]{biblatex}
\addbibresource{ref.bib} %BibTeX数据文件及位置
\setbeamerfont{footnote}{size=\tiny} %调整注脚的文字大小

% other packages
\usepackage{latexsym,amsmath,xcolor,multicol,booktabs,calligra}
\usepackage{graphicx,pstricks,listings,stackengine}
\usepackage{pgfplots}
\pgfplotsset{width=\textwidth,compat=1.9}
\setCJKfamilyfont{myfont}{msyh.ttc}
\newcommand{\MyFont}{\CJKfamily{myfont}}
%\usepackage{auto-pst-pdf}
%\usepackage[utf8]{inputenc}

\author{\MyFont{Leyi Zhao}}
\title{\MyFont{Separating Out the Eigenvalue Densities: Computing the Jacobians}}
\subtitle{Random Matrix Theory with Its Applications}
\institute{\MyFont{School of Cyberspace}}
\date{\today}
\usepackage{hdu}

% defs
\def\cmd#1{\texttt{\color{red}\footnotesize $\backslash$#1}}
\def\env#1{\texttt{\color{blue}\footnotesize #1}}
\definecolor{deepblue}{rgb}{0,0,0.5}
\definecolor{deepred}{rgb}{0.6,0,0}
\definecolor{deepgreen}{rgb}{0,0.5,0}
\definecolor{halfgray}{gray}{0.55}

\lstset{
    basicstyle=\ttfamily\small,
    keywordstyle=\bfseries\color{deepblue},
    emphstyle=\ttfamily\color{deepred},    % Custom highlighting style
    stringstyle=\color{deepgreen},
    numbers=left,
    numberstyle=\small\color{halfgray},
    rulesepcolor=\color{red!20!green!20!blue!20},
    frame=shadowbox,
}

%%tips:
\usepackage{geometry}
\usepackage{lipsum}
\usepackage[dvipsnames,svgnames]{xcolor}
\usepackage{tcolorbox}
\tcbuselibrary{most}
\definecolor{tipscolor}{rgb}{0.77,0.72,0.65} % 莫兰迪棕色
\definecolor{OldLace}{rgb}{0.992156863, 0.9607843, 0.9019608}
% ------------------******-------------------
\newtcolorbox{tips}[2][]
{enhanced,breakable,
left=12pt,right=12pt,% 左右边距
% fonttitle=\bfseries, % 可以设置标题是否粗体
coltitle=white, % 标题字体颜色
colbacktitle=tipscolor, % 标题背景颜色
attach boxed title to top left={yshifttext=-1mm},
boxed title style={skin=enhancedfirst jigsaw,arc=1mm,bottom=0mm,boxrule=0mm},
boxrule=1pt, % 边框线宽
colback=OldLace, % 文本框背景颜色
colframe=tipscolor, % 框线颜色
sharp corners=northwest,
% drop fuzzy shadow, % 可以选择是否设置阴影效果
title=\vspace{3mm}#2,
arc=1mm,
#1}
%%


\begin{document}

%\kaishu
\MyFont
\begin{frame}
    \titlepage
    \begin{figure}[htpb]
        \begin{center}
            \includegraphics[width=0.18\linewidth]{img/hdu-logo.jpg} %插入学校的logo，支持多种格式
        \end{center}
    \end{figure}
\end{frame}

%  

\begin{frame}
    \tableofcontents[sectionstyle=show,subsectionstyle=show/shaded/hide,subsubsectionstyle=show/shaded/hide]
\end{frame}

\section{Background}
\begin{frame}{Spectral Theorem}
    \begin{itemize}
        \small
        \item The spectral theorem:
              \begin{align*}
                  M & = U \Lambda U^*  (unitary, symplectic) \\
                  M & = U \Lambda U^T  (orthogonal)
              \end{align*}
              $M$ can be any kinds of matrices.
        \item Three ensembles:\footfullcite{yin2014convergence}
              \begin{itemize}
                  \small
                  \item \textbf{Gaussian Orthogonal Ensemble}: the set of $N \times N$ random \textcolor{red}{real} \textbf{symmetric} matrices
                  \item \textbf{Gaussian Unitary Ensemble}: the set of $N \times N$ random \textcolor{red}{complex} \textbf{Hermitian} matrices
                  \item \textbf{Gaussian Symplectic Ensemble}: the set of $N \times N$ random \textcolor{red}{quaternion}(四元数) \textbf{self-dual Hermitian} matrices
              \end{itemize}
        \item 这可以被视为一种坐标系变化：$M \mapsto (\Lambda, U)$
        \item 此时可以考虑计算雅可比行列式（Jacobian）：$\left|det \frac{\partial M}{\partial\left(\Lambda, U\right)}\right|$
    \end{itemize}
\end{frame}

\begin{frame}{Spectral Theorem}
    \begin{itemize}
        \item Every \textcolor{red}{real symmetric} matrix $S$ with \textcolor{red}{real} eigenvalues in $\Lambda$ and orthonormal eigenvectors in the columns of $Q$(\textcolor{red}{orthogonal}(正交)) can be diagonalized:\footfullcite{strang2016introduction}
              \begin{align*}
                  S=Q\Lambda Q^{-1} = Q\Lambda Q^T \quad\text{with}\quad Q^{-1} = Q^T
              \end{align*}
        \item Every \textcolor{red}{Hermitian} matrix $S$ with \textcolor{red}{real} eigenvalues in $\Lambda$ and orthonormal eigenvectors in the columns of $U$(\textcolor{red}{unitary}(酉)) can be diagonalized:\footfullcite{strang2016introduction}
              \begin{align*}
                  S=U\Lambda U^{-1} = U\Lambda U^* \quad\text{with}\quad U^{-1} = U^*
              \end{align*}
    \end{itemize}
\end{frame}

\begin{frame}{Gaussian Orthogonal Ensemble(GOE)}
    \begin{itemize}
        \item \textbf{Gaussian Orthogonal Ensemble}: the set of $N \times N$ random \textcolor{red}{real} \textbf{symmetric} matrices
        \item Symmetric matrix:
              \begin{itemize}
                  \item $S=S^T$
                  \item $S=X\Lambda X^{-1}, S^T=(X^{-1})^T\Lambda X^T\quad \text{when}\quad S=S^T, X^TX=I$
                  \item A symmetric matrix has only \emph{real eigenvalues}.
                  \item Symmetric diagonalization: $S=Q\Lambda Q^{-1} = Q\Lambda Q^T \quad\text{with}\quad Q^{-1} = Q^T$
                  \item Example: $\begin{bmatrix}
                                a & b \\
                                b & c
                            \end{bmatrix}$
                  \item Free variables number is $n+\frac{n(n-1)}{2}\cdot 1$
              \end{itemize}
    \end{itemize}
\end{frame}

\begin{frame}{Gaussian Unitary Ensemble(GUE)}
    \begin{itemize}
        \item \textbf{Gaussian Unitary Ensemble}: the set of $N \times N$ random \textcolor{red}{complex} \textbf{Hermitian} matrices
        \item Hermitian matrices:
              \begin{itemize}
                  \item $S=S^*$
                  \item $S=U\Lambda U^{-1}, S^*=(U^{-1})^*\Lambda U^*\quad \text{when}\quad S=S^*, U^*U=I$
                  \item Hermitian diagonalization: $S=U\Lambda U^{-1} = U\Lambda U^* \quad\text{with}\quad U^{-1} = U^*$
                  \item Example: $\begin{bmatrix}
                                a    & b+ci \\
                                b-ci & d
                            \end{bmatrix}$
                  \item Free variables number is $n+\frac{n(n-1)}{2}\cdot 2$
              \end{itemize}
    \end{itemize}
\end{frame}

\begin{frame}{Gaussian Symplectic Ensemble(GSE)}
    \begin{itemize}
        \item \textbf{Gaussian Symplectic Ensemble}: the set of $N \times N$ random \textcolor{red}{quaternion} \textbf{self-dual Hermitian} matrices\footfullcite{许方官2012四元数物理学} \footfullcite{yin2014convergence}
        \item Quaternions were invented in 1843 by the Irish mathematician
              William Rowen Hamilton as an extension of complex numbers into three dimensions\footfullcite{deavours1973quaternion} \footfullcite{hamilton1866elements}, and it's well known that the quaternion field $\mathbb{Q} $ can be represented as a two-dimensional
              complex vector space\footfullcite{chevalley1946theory}.
    \end{itemize}
\end{frame}

\begin{frame}{四元数}
    \begin{itemize}
        \item 一个三位空间的虚矢量$\mathbf{a}=a_1\mathbf{i}+a_2\mathbf{j}+a_3\mathbf{k}$与一个实数$a_0$组合成一个数$A=a_0+\mathbf{a}$被称为\textbf{四元数}。其中实数$a_0$被称为四元数$A$的\textbf{标部}，实矢量$\mathbf{a}$称为四元数$A$的矢部。
        \item 在四元数$A$的矢部中，三个矢量$\mathbf{i},\mathbf{j},\mathbf{k}$是三维空间中三个互相垂直且方向固定的单位矢量。
        \item 一个四元数实际上是由四个基$1,\mathbf{i},\mathbf{j},\mathbf{k}$的线性组合构成的。
    \end{itemize}
\end{frame}

\begin{frame}{四元数代数运算规则}
    \begin{itemize}
        \item \textbf{相等}: 两个四元数$A$和$B$，若他们相应的四个分量分别相等，即
              \begin{align*}
                  a_n=b_n, \quad n=0,1,2,3
              \end{align*}
              则称他们\textbf{相等}，记作
              \begin{align*}
                  A=B
              \end{align*}
              一个四元数的等式相当于四个实数等式。
        \item \textbf{加法}:两个四元数$A$和$B$之和$C$仍为四元数，记作
              \begin{align*}
                  C=A+B
              \end{align*}
              其定义为
              \begin{align*}
                  c_n=a_n+b_n,\quad n=0,1,2,3
              \end{align*}
              由实数加法的交换律、结合律知，四元数加法也存在着交换律和结合律。
    \end{itemize}
\end{frame}

\begin{frame}{四元数代数运算规则}
    \begin{itemize}
        \small
        \item \textbf{乘法}:两个四元数$A$和$B$乘积的定义是
              \begin{align*}
                  \begin{split}
                      AB & =\left(a_0+\mathbf{a}\right)\left(b_0+\mathbf{b}\right)                                  \\
                         & =a_0b_0-\mathbf{a}\cdot\mathbf{b}+a_0\mathbf{b}+b_0\mathbf{a}+\mathbf{a}\times\mathbf{b} \\
                      BA & =\left(b_0+\mathbf{b}\right)\left(a_0+\mathbf{a}\right)                                  \\
                         & =a_0b_0-\mathbf{a}\cdot\mathbf{b}+a_0\mathbf{b}+b_0\mathbf{a}+\mathbf{b}\times\mathbf{a}
                  \end{split}
              \end{align*}
              两个四元数的乘积仍是四元数，但因为$\mathbf{a}\times\mathbf{b}\neq\mathbf{b}\times\mathbf{a}$，故$AB\neq BA$

              四元数的乘法规则是由三个互相垂直的空间单位矢量$\mathbf{i},\mathbf{j},\mathbf{k}$之间的乘法规则
              \begin{align*}
                  \left\{\begin{matrix}\begin{aligned}
                          \mathbf{ii} & =\mathbf{jj}=\mathbf{kk}= \mathbf{ijk} = -1 \\
                          \mathbf{ij} & =-\mathbf{ji}=\mathbf{k}                    \\
                          \mathbf{jk} & =-\mathbf{kj}=\mathbf{i}                    \\
                          \mathbf{ki} & =-\mathbf{ik}=\mathbf{j}
                      \end{aligned}
                         \end{matrix}\right.
              \end{align*}
              所导致的
    \end{itemize}
\end{frame}

\begin{frame}{四元数代数运算规则}
    \begin{itemize}
        \scriptsize
        \item \textbf{共轭(四元共轭)}:与四元数$A=a_0+\mathbf{a}$对应的另一个四元数
              \begin{align*}
                  \overline{A} = a_0 - \mathbf{a}
              \end{align*}
              称为A的\textbf{四元共轭}，反之，A也称为$\overline{A}$的\textbf{四元共轭}

              此外，由
              \begin{align*}
                  \begin{split}
                      \overline{AB}                      & =a_0b_0-\mathbf{a}\cdot\mathbf{b}-a_0\mathbf{b}-b_0\mathbf{a}-\mathbf{a}\times\mathbf{b} \\
                      \overline{B}\thinspace\overline{A} & =\left(b_0-\mathbf{b}\right)\left(a_0-\mathbf{a}\right)                                  \\
                                                         & =a_0b_0-\mathbf{a}\cdot\mathbf{b}-a_0\mathbf{b}-b_0\mathbf{a}+\mathbf{b}\times\mathbf{a}
                  \end{split}
              \end{align*}
              可知
              \begin{align*}
                  \overline{AB}=\overline{B}\thinspace\overline{A}
              \end{align*}
        \item \textbf{模方}:称实数
              \begin{align*}
                  \left \| A \right \| =\left \| \overline{A} \right \| =A\overline{A}=\overline{A}A=a_0^2+\mathbf{a}\cdot\mathbf{a}
              \end{align*}
              为四元数$A$或$\overline{A}$的\textbf{模方}，且两个四元数之积的模方等于两个四元数模方之积$\left \| AB \right \|=\left \| A \right \|\left \| B \right \|$
    \end{itemize}
\end{frame}

\begin{frame}{四元数的矩阵表示}
    \begin{itemize}
        \scriptsize
        \item \textbf{一级四元数与$2\times 2$矩阵}: 把四元数的四个基$1,\mathbf{i},\mathbf{j},\mathbf{k}$分别用四个$2\times 2$矩阵表示成
              \begin{align*}
                  1=\begin{bmatrix} 1 & 0\\ 0 & 1\end{bmatrix}, \quad\mathbf{i}=\begin{bmatrix} i & 0\\ 0 & -i\end{bmatrix},\quad\mathbf{j}=\begin{bmatrix} 0 & 1\\ -1 & 0\end{bmatrix},\quad\mathbf{k}=\begin{bmatrix} 0 & i\\ i & 0\end{bmatrix}
              \end{align*}
              这样，一个四元数可以用一个$2\times 2$的复矩阵表示成
              \begin{align*}
                  A=a_0+\mathbf{a}=\begin{bmatrix} a_0+ia_1 & a_2+ia_3\\ -a_2+ia_3 & a_0-ia_1\end{bmatrix}=\begin{bmatrix} \alpha & \beta\\ -\bar{\beta} & \bar{\alpha}\end{bmatrix}
              \end{align*}
        \item \textbf{共轭}:
              \begin{align*}
                  \overline{A}=a_0-a_1\mathbf{i}-a_2\mathbf{j}-a_3\mathbf{k}=\begin{bmatrix} a_0-ia_1 & -a_2-ia_3\\ a_2-ia_3 & a_0+ia_1\end{bmatrix}=\begin{bmatrix} \bar{\alpha} & -\beta\\ \bar{\beta} & \alpha\end{bmatrix}
              \end{align*}
        \item \textbf{模方}:四元数矩阵的模方就是其表示的矩阵的行列式，即
              \begin{align*}
                  \left\|A\right\|=\det A=\sqrt{a_0^2+a_1^2+a_2^2+a_3^2}=\sqrt{\left|\alpha\right|^2+\left|\beta\right|^2}
              \end{align*}

    \end{itemize}
\end{frame}

\begin{frame}{四元数的矩阵表示}
    \tiny
    一个四元数可以用$2\times 2$矩阵表示，类似地，二级四元数可以用$4\times 4$矩阵表示，其四个基$1,\mathbf{i_2},\mathbf{j_2},\mathbf{k_2}$表示成
    \begin{align*}
        1=\begin{bmatrix} I_2 & 0\\ 0 & I_2\end{bmatrix},\quad\mathbf{i}=\begin{bmatrix} iI_2 & 0\\ 0 & -iI_2\end{bmatrix},\quad\mathbf{j}=\begin{bmatrix} 0 & I_2\\ -I_2 & 0\end{bmatrix},\quad\mathbf{k}=\begin{bmatrix} 0 & iI_2\\ iI_2 & 0\end{bmatrix}
    \end{align*}
    其中$I_2$是$2\times 2$的单位矩阵。这样，一个二级四元数的矩阵表示为
    \begin{align*}
        \begin{split}
            Q= & A+B\mathbf{i_2}+C\mathbf{j_2}+D\mathbf{k_2}                                                                                                                                                                                                                                                                                                                                                                                                                                                                                                                       \\
            =  & \begin{bmatrix} A+iB & C+iD\\ -C+iD & A-iB\end{bmatrix}                                                                                                                                                                                                                                                                                                                                                                                                                                                                                                           \\
            =  & \begin{bmatrix}a_0-b_1+\mathrm{i}(a_1+b_0)\!&\!a_2-b_3+\mathrm{i}(a_3+b_2)\!&\!c_0-d_1+\mathrm{i}(c_1+d_0)\!&\!c_2-d_3+\mathrm{i}(c_3+d_2)\\-(a_2+b_3)+\mathrm{i}(a_3-b_2)\!&\!a_0+b_1-\mathrm{i}(a_1-b_0)\!&\!-(c_2+d_3)+\mathrm{i}(c_3-d_2)\!&\!c_0+d_1-\mathrm{i}(c_1-d_0)\\-(c_0+d_1)-\mathrm{i}(c_1-d_0)\!&\!-(c_2+d_3)-\mathrm{i}(c_3-d_2)\!&\!a_0+b_1+\mathrm{i}(a_1-b_0)\!&\!a_2+b_3+\mathrm{i}(a_3-b_2)\\c_2-d_3-\mathrm{i}(c_3+d_2)\!&\!-(c_0-d_1)+\mathrm{i}(c_1+d_0)\!&\!-(a_2-b_3)+\mathrm{i}(a_3+b_2)\!&\!a_0-b_1-\mathrm{i}(a_1+b_0)\end{bmatrix}.
        \end{split}
    \end{align*}
\end{frame}

\begin{frame}{Gaussian Symplectic Ensemble(GSE)}
    \begin{itemize}
        \scriptsize
        \item An $n\times n$ \textbf{quaternion self-dual Hermitian matrix} $^q\mathbf{A}_n=(x_{jk})_{n\times n}$ is a matrix whose entries $x_j(j,k=1,\ldots,n)$ are \textbf{quaternions} and satisfy $x_{jk}=\bar{x}_{kj}.$ If each $x_{jk}$ is viewed as a $2\times2$ matrix, then $\mathbf{A}_n=(x_{jk})$ is in fact a $2n\times2n$ Hermitian matrix. And, the entries of $\mathbf{A}_n$ can be represented as
              \begin{align*}
                  x_{jk}=\begin{bmatrix}a_{jk}+b_{jk}i&c_{jk}+d_{jk}i\\-c_{jk}+d_{jk}i&a_{jk}-b_{jk}i\end{bmatrix}=\begin{bmatrix}\alpha_{jk}&\beta_{jk}\\-\overline{\beta}_{jk}&\overline{\alpha}_{jk}\end{bmatrix},1\leq j<k\leq n,
              \end{align*}
              and $x_{jj}=\begin{bmatrix}a_{jj}&0\\0&a_{jj}\end{bmatrix}$ , where $a_{jk},b_{jk},c_{jk},d_{jk}\in\mathbb{R}$ and $1\leq j,k\leq n$. It is well known that the multiplicities of all the eigenvalues of $\mathcal{A}_n$ are even\footfullcite{zhang1997quaternions}.
        \item Example: $\begin{bmatrix}
                      \begin{bmatrix}a_{11}&0\\0&a_{11}\end{bmatrix}                                            & \begin{bmatrix}a_{12}+b_{12}i&c_{12}+d_{12}i\\-c_{12}+d_{12}i&a_{12}-b_{12}i\end{bmatrix} \\
                      \begin{bmatrix}a_{12}-b_{12}i&-c_{12}-d_{12}i\\c_{12}-d_{12}i&a_{12}+b_{12}i\end{bmatrix} & \begin{bmatrix}a_{44}&0\\0&a_{44}\end{bmatrix}
                  \end{bmatrix}$
        \item Free variables number is $n+\frac{n(n-1)}{2}\cdot 4$
    \end{itemize}
\end{frame}

\section{A Two-Dimensional Example}

\begin{frame}{$2\times 2$ Real Symmetric Matrix}
    \begin{itemize}
        \scriptsize
        \item Consider a $2\times 2$ real symmetric matrix $M=\begin{bmatrix}a&b\\b&c\end{bmatrix}$ with distinct eigenvalues $\lambda_1\neq\lambda_2$. Then $M=U\begin{bmatrix}\lambda_1&0\\0&\lambda_2\end{bmatrix}U^T$, and without loss of generality we can assume $U=\begin{bmatrix} \cos\theta &-\sin\theta \\ \sin\theta & \cos\theta \end{bmatrix}$ for some $0\leq\theta<2\pi$.
        \item $U$是二维空间的一组基构成的矩阵，$\begin{bmatrix}\cos\theta \\ \sin\theta\end{bmatrix}$与$\begin{bmatrix}-\sin\theta \\ \cos\theta\end{bmatrix}$相互正交
        \item
              \begin{align*}
                  \begin{split}
                      a & =\lambda_{1}\cos^{2}\theta+\lambda_{2}\sin^{2}\theta \\
                      b & =(\lambda_1-\lambda_2)\cos\theta\sin\theta           \\
                      c & =\lambda_{1}\sin^{2}\theta+\lambda_{2}\cos^{2}\theta
                  \end{split}
              \end{align*}
        \item
              \begin{align*}
                  \frac{\partial(a,b,c)}{\partial(\lambda_1,\lambda_2,\theta)}=\begin{bmatrix}\cos^2\theta&\sin^2\theta&-2(\lambda_1-\lambda_2)\cos\theta\sin\theta\\\cos\theta\sin\theta&-\cos\theta\sin\theta&(\lambda_1-\lambda_2)(\cos^2\theta-\sin^2\theta)\\\sin^2\theta&\cos^2\theta&2(\lambda_1-\lambda_2)\cos\theta\sin\theta\end{bmatrix}.
              \end{align*}
    \end{itemize}
\end{frame}

\begin{frame}{$2\times 2$ Real Symmetric Matrix}
    \begin{itemize}
        \scriptsize
        \item Since
              \begin{align*}
                  \frac{\partial(a,b,c)}{\partial(\lambda_1,\lambda_2,\theta)}=\begin{bmatrix}\cos^2\theta&\sin^2\theta&-2(\lambda_1-\lambda_2)\cos\theta\sin\theta\\\cos\theta\sin\theta&-\cos\theta\sin\theta&(\lambda_1-\lambda_2)(\cos^2\theta-\sin^2\theta)\\\sin^2\theta&\cos^2\theta&2(\lambda_1-\lambda_2)\cos\theta\sin\theta\end{bmatrix}.
              \end{align*}
        \item The determinant can be:
              \begin{align*}
                  \left|\det\frac{\partial(a,b,c)}{\partial(\lambda_1,\lambda_2,\theta)}\right|=|\lambda_1-\lambda_2|f_1(\theta)
              \end{align*}
              where
              \begin{align*}
                  \left.f_1(\theta)=\left|\det\begin{bmatrix}\cos^2\theta&\sin^2\theta&-2\sin2\theta\\\frac12\sin2\theta&-\frac12\sin2\theta&\cos2\theta\\\sin^2\theta&\cos^2\theta&\sin2\theta\end{bmatrix}\right.\right|=1>0
              \end{align*}
    \end{itemize}
\end{frame}

\begin{frame}{$2\times 2$ Hermitian Matrix}
    \begin{itemize}
        \tiny
        \item Consider a $2\times 2$ Hermitian matrix $M=\begin{bmatrix}a & b+ci \\b-ci & d\end{bmatrix}$with distinct eigenvalues $\lambda_1\neq\lambda_2$. Then $M=U\begin{bmatrix}\lambda_1&0\\0&\lambda_2\end{bmatrix}U^*$, and without loss of generality we can assume $U=\begin{bmatrix}\cos\theta & -\left(\cos\chi-i\sin\chi\right)\sin\theta\\ \left(\cos\chi+i\sin\chi\right)\sin\theta & \cos\theta\end{bmatrix}$ for some $0\leq\theta\leq\frac12\pi$ and $\quad0\leq\chi<2\pi$.\footfullcite{howard2012diagonalization}
        \item \begin{align*}
                  \begin{split}
                      a & =\lambda_1\cos^2\theta+\lambda_2\sin^2\theta                   \\
                      b & =\left(\lambda_1-\lambda_2\right)\cos\theta\sin\theta\cos\chi  \\
                      c & =\left(-\lambda_1+\lambda_2\right)\cos\theta\sin\theta\sin\chi \\
                      d & =\lambda_{1}\sin^{2}\theta+\lambda_{2}\cos^{2}\theta
                  \end{split}
              \end{align*}
        \item
              \begin{align*}
                  \begin{split}
                       & \frac{\partial(a,b,c,d)}{\partial(\lambda_1,\lambda_2,\theta,\chi)}=                                                                                                                                                                                                                                                                                                                                                                                                                                                                                                                                                                                                                          \\
                       & \begin{bmatrix}\cos^2\theta \!&\! \sin^2\theta \!&\! -2(\lambda_1\!-\!\lambda_2)\cos\theta\!\sin\theta \!&\! 0\\ \cos\theta\!\sin\theta\!\cos\chi \!&\! -\!\cos\theta\!\sin\theta\!\cos\chi \!&\! (\lambda_1\!-\!\lambda_2)(\cos^2\theta\!-\!\sin^2\theta)\cos\chi \!&\! -\left(\lambda_1\!-\!\lambda_2\right)\sin\theta\!\cos\theta\!\sin\chi \\ -\cos\theta\!\sin\theta\!\sin\chi \!&\! \cos\theta\!\sin\theta\!\sin\chi \!&\! -(\lambda_1\!-\!\lambda_2)(\cos^2\theta\!-\!\sin^2\theta)\sin\chi \!&\! -\left(\lambda_1\!-\!\lambda_2\right)\sin\theta\!\cos\theta\!\cos\chi \\ \sin^2\theta \!&\! \cos^2\theta \!&\! 2(\lambda_1\!-\!\lambda_2)\cos\theta\!\sin\theta \!&\! 0\end{bmatrix}
                  \end{split}
              \end{align*}
    \end{itemize}
\end{frame}

\begin{frame}{$2\times 2$ Hermitian Matrix}
    \begin{itemize}
        \tiny
        \item Since
              \begin{align*}
                  \begin{split}
                       & \frac{\partial(a,b,c,d)}{\partial(\lambda_1,\lambda_2,\theta,\chi)}=                                                                                                                                                                                                                                                                                                                                                                                                                                                                                                                                                                                                                          \\
                       & \begin{bmatrix}\cos^2\theta \!&\! \sin^2\theta \!&\! -2(\lambda_1\!-\!\lambda_2)\cos\theta\!\sin\theta \!&\! 0\\ \cos\theta\!\sin\theta\!\cos\chi \!&\! -\!\cos\theta\!\sin\theta\!\cos\chi \!&\! (\lambda_1\!-\!\lambda_2)(\cos^2\theta\!-\!\sin^2\theta)\cos\chi \!&\! -\left(\lambda_1\!-\!\lambda_2\right)\sin\theta\!\cos\theta\!\sin\chi \\ -\cos\theta\!\sin\theta\!\sin\chi \!&\! \cos\theta\!\sin\theta\!\sin\chi \!&\! -(\lambda_1\!-\!\lambda_2)(\cos^2\theta\!-\!\sin^2\theta)\sin\chi \!&\! -\left(\lambda_1\!-\!\lambda_2\right)\sin\theta\!\cos\theta\!\cos\chi \\ \sin^2\theta \!&\! \cos^2\theta \!&\! 2(\lambda_1\!-\!\lambda_2)\cos\theta\!\sin\theta \!&\! 0\end{bmatrix}
                  \end{split}
              \end{align*}
        \item The determinant can be:
              \begin{align*}
                  \left|\det\frac{\partial(a,b,c,d)}{\partial(\lambda_1,\lambda_2,\theta,\chi)}\right|=\left(\lambda_1-\lambda_2\right)^2f_2(\theta, \chi)
              \end{align*}
              where
              \begin{align*}
                  \left.f_2(\theta,\chi)=\left|\det\begin{bmatrix}\cos^2\theta & \sin^2\theta & -\sin2\theta & 0\\ \cos\theta\sin\theta\cos\chi & -\cos\theta\sin\theta\cos\chi & \cos2\theta\cos\chi & -\frac{1}{2}\sin2\theta\sin\chi \\ -\cos\theta\sin\theta\sin\chi & \cos\theta\sin\theta\sin\chi & -\cos2\theta\sin\chi & -\frac{1}{2}\sin2\theta\cos\chi \\ \sin^2\theta & \cos^2\theta & \sin2\theta & 0\end{bmatrix}\right.\right|
              \end{align*}
              is independent of the eigenvalues of M.
    \end{itemize}
\end{frame}

\begin{frame}{$4 \times 4$ Hermitian self-dual matrix}
    \begin{itemize}
        \scriptsize
        \item Consider a $4\times 4$ Hermitian self-dual matrix
              $M=\begin{bmatrix}
                      \begin{bmatrix}a_{11}&0\\0&a_{11}\end{bmatrix}                                            & \begin{bmatrix}a_{12}+b_{12}i&c_{12}+d_{12}i\\-c_{12}+d_{12}i&a_{12}-b_{12}i\end{bmatrix} \\
                      \begin{bmatrix}a_{12}-b_{12}i&-c_{12}-d_{12}i\\c_{12}-d_{12}i&a_{12}+b_{12}i\end{bmatrix} & \begin{bmatrix}a_{44}&0\\0&a_{44}\end{bmatrix}
                  \end{bmatrix}$. Then $M=U\begin{bmatrix}\mu_1&0&0&0\\0&\mu_1&0&0\\0&0&\mu_2&0\\0&0&0&\mu_2\end{bmatrix}U^*$.
        \item The determinant can be:
              \begin{align*}
                  \left|\det\frac{\partial(a_{11},a_{12},a_{44},b_{12},c_{12},d_{12})}{\partial(\Lambda, U)}\right|=\left(\lambda_1-\lambda_2\right)^4f_4(U)
              \end{align*}
    \end{itemize}
\end{frame}

\begin{frame}{Conclusion}
    \begin{itemize}
        \scriptsize
        \item A well-known analogy between RMT and statistical machanics $\beta$ corresponds to an inverse temperature.
              \begin{align*}
                  \text{Free veriables: }n+\frac{n(n-1)}{2}\cdot \beta \\
                  \left\{\begin{matrix}\begin{aligned}
                          \beta=1 & \quad\text{orthogonal ensembles,} \\
                          \beta=2 & \quad\text{unitary ensembles,}    \\
                          \beta=4 & \quad\text{symplectic ensembles.}
                      \end{aligned}
                         \end{matrix}\right.
              \end{align*}
        \item Similarly, as for $2\times 2$ matrix, the determinant is:
              \begin{align*}
                  \left|\det\frac{\partial(M)}{\partial(\Lambda, U)}\right|=\left(\lambda_1-\lambda_2\right)^{\beta}f_{\beta}(U) \\
                  \left\{\begin{matrix}\begin{aligned}
                          \beta=1 & \quad\text{orthogonal ensembles,} \\
                          \beta=2 & \quad\text{unitary ensembles,}    \\
                          \beta=4 & \quad\text{symplectic ensembles.}
                      \end{aligned}
                         \end{matrix}\right.
              \end{align*}
    \end{itemize}
\end{frame}

\begin{frame}
    \begin{center}
        {\Huge\calligra Thanks!}
    \end{center}
\end{frame}

\end{document}